%+------------------------------------------------------------------------+
%| NOTE: Shared academic application draft; no submission authorised.
%| DATE: 2026-10-23
%| Author: Antonio Michele Miti
%| Tools: OpenAI ChatGPT/Codex; historical model versions not reconstructed.
%| Credits: Cv/data/cv.yaml; Cv/data/publications.bib;
%|          Teaching-Statement/teaching-statement-Miti.tex; primary metadata.
%+------------------------------------------------------------------------+
% FirstPrinciples adaptation: research profile reframed; factual records preserved.
% New AI narrative is red. Source-backed factual records are not a claim of
% independent personal verification. Scheduled unconfirmed stays and dummy
% supervision entries are excluded. Pending author review; not submitted.
\documentclass[11pt,a4paper]{article}
\usepackage[a4paper,margin=2.2cm]{geometry}
\usepackage[T1]{fontenc}
\usepackage[utf8]{inputenc}
\usepackage[english]{babel}
\usepackage{lmodern,microtype,amsmath,amssymb,xcolor,enumitem,fancyhdr}
\usepackage[hidelinks]{hyperref}
\setlength{\parindent}{0pt}
\setlength{\parskip}{0.35em}
\setlist[itemize]{leftmargin=1.2em,itemsep=0.15em,topsep=0.2em}
\setlist[enumerate]{leftmargin=1.8em,itemsep=0.5em,topsep=0.2em}
\newcommand{\AI}[1]{{\color{red}#1}}
\pagestyle{fancy}\fancyhf{}\renewcommand{\headrulewidth}{0pt}
\fancyfoot[L]{\small Antonio Michele Miti | Application draft}
\fancyfoot[R]{\small\thepage}
\begin{document}
\begin{center}
{\LARGE\bfseries Antonio Michele Miti}\\[0.2em]
{\large Academic curriculum vitae -- draft for review}\\[0.3em]
\href{mailto:antoniomichele.miti@gmail.com}{antoniomichele.miti@gmail.com}\quad
\href{https://www.antoniomiti.it/}{antoniomiti.it}\\
\href{https://orcid.org/0000-0002-8829-1943}{ORCID: 0000-0002-8829-1943}\quad 23 October 2026
\end{center}
\section*{Research}
\AI{Mathematical physics and differential geometry: multisymplectic observables, constraint reduction and higher structures, with theoretical-physics training and joint work on quantum-circuit uncomputation.}
% Employment end date confirmed by the author in COMUNI-CV-01, 23 October 2026.
\section*{Research appointments}
\begin{itemize}
\item \textbf{February 2024--April 2026:} CIVIS-3i postdoctoral fellow, Department of Mathematics, Sapienza Universit\`a di Roma. MSCA-H2020 COFUND fellowship on multisymplectic geometry and classical field theories.
\item \textbf{October 2022--January 2024:} Postdoctoral fellow, Universit\`a Cattolica del Sacro Cuore, Brescia. Symplectic and multisymplectic geometry with applications.
\item \textbf{October 2021--September 2022:} Postdoctoral fellow, Max Planck Institute for Mathematics, Bonn. Higher structures and multisymplectic geometry.
\end{itemize}
\section*{Education}
\begin{itemize}
\item \textbf{1 April 2021:} PhD in Mathematics, joint degree of Universit\`a Cattolica del Sacro Cuore and KU Leuven. Thesis: \emph{Homotopy comomentum maps in multisymplectic geometry}. Supervisors: Mauro Spera and Marco Zambon.
\item \textbf{November 2015:} Master's degree in Theoretical Physics, Universit\`a degli Studi di Milano, 110/110. Thesis on algebraic quantization of Jacobi fields and the geometric approach to Peierls brackets. Supervisors: Claudio Dappiaggi and Livio Pizzocchero.
\item \textbf{October 2010:} Bachelor's degree in Physics, Universit\`a degli Studi di Milano-Bicocca. Thesis on Lie groups and rigid body mechanics. Supervisor: Franco Magri.
\end{itemize}
\section*{Teaching and graduate learning activities}
\begin{itemize}
\item \textbf{2024, Salerno:} Co-lecturer with Leonid Ryvkin, research minicourse \emph{Reduction of Multisymplectic Observables}. Eight-hour joint course (total course duration).
\item \textbf{2017--2018, 2022 and 2023, Brescia:} Teaching assistant, \emph{Elements of Higher Geometry}, course led by Mauro Spera. Exercises and selected lectures; support for assessment.
\item \textbf{2018, Brescia:} Teaching assistant, \emph{Differential Topology and Geometry}. Exercises in differential topology and simplicial homology; support for assessment.
\item \textbf{2014--2015, Milano:} Student support at the Physics Department's Computation and Multimedia Laboratory. Computing environments and scientific programming projects, including C, Python and CUDA.
\item \textbf{2024, Lyon:} Organiser, reading seminar \emph{Diffeological Structures in Lagrangian Field Theory}.
\item \textbf{2025, Sapienza:} Co-organiser and host of the doctoral course on perturbative algebraic quantum field theory taught by Fabrizio Zanello.
\end{itemize}
\section*{Selected research visits}
\begin{itemize}
\item \textbf{March 2025:} Institut Mittag-Leffler, programme \emph{Cohomological Aspects of Quantum Field Theory}.
\item \textbf{October--December 2024:} Institut Camille Jordan, Universit\'e Claude Bernard Lyon 1, CIVIS-3i secondment.
\item \textbf{December 2023--January 2024:} Institut Camille Jordan, supported by an INdAM research-abroad grant.
\end{itemize}
\section*{Selected grants and distinctions}
\begin{itemize}
\item \textbf{2026:} MSCA Seal of Excellence for the GROCOMO-CFT proposal (MSCA-PF-2025). \AI{This is a quality distinction, not a claim that the proposal was funded.}
\item \textbf{2026:} Sapienza \emph{Avvio alla ricerca 2025}, EUR 2,800, supporting the GroCO research project.
\item \textbf{2023:} INdAM research-abroad grant, EUR 3,000.
\end{itemize}
\section*{Academic service and organisation}
\begin{itemize}
\item Refereeing recorded since 2023 for SIGMA and Geometriae Dedicata.
\item Co-organisation of workshops on multisymplectic geometry: KU Leuven (2020) and Lyon (2024).
\item Assistance with departmental computer systems, Universit\`a Cattolica del Sacro Cuore (2017--2018).
\end{itemize}
\section*{Languages and technical background}
Italian: native. English: TOEFL iBT 102/120, April 2018 (historical test result). French: beginner, as recorded in the source CV. Scientific computing and document preparation: Python, C/C++, CUDA, \LaTeX\ and Linux.

\end{document}
